% 分章页60—65，2.7与Software Package；印刷页132—137。
\section{模拟与真实数据分析}
\subsection{模拟}

通过表型网络的模拟研究，评价基于评分的网络结构推断方法的性能。模拟网络包含30个节点和190条有向边。使用稀疏结构方程模型建立表型网络，并计算节点评分度量。用于随机生成网络结构及其参数的结构方程模型为
\begin{equation}
Y=\epsilon\Gamma^{-1},\tag{2.166}
\end{equation}
其中 $\epsilon\sim N(0,0.05\times I)$。对于表型子网络，若存在从节点 $j$ 指向节点 $i$ 的边，则建模所用结构方程中的参数 $\Gamma_{ij}$ 从区间 $(0.5,1)$ 或 $(-1,-0.5)$ 上的均匀分布随机生成；否则令 $\Gamma_{ij}=0$。

随机生成的表型网络中，每个节点的期望度数为3。模拟重复5000次。用五折交叉验证确定惩罚参数 $\lambda$，并在功效模拟中用它推断网络。使用检出功效（power of detection，PD）与错误发现率（FDR）两个指标，评价算法识别网络结构的性能。具体而言，令 $N_t$ 为5000次网络重复中的边总数，$\widehat N_t$ 为推断算法检出的边总数，$N_{\mathrm{True}}$ 为模拟网络中检出的真边总数，$N_{\mathrm{False}}$ 为 $\widehat N_t$ 中检出的假边数。检出功效定义为 $N_{\mathrm{True}}/N_t$，错误发现率定义为 $N_{\mathrm{False}}/\widehat N_t$。

比较五种统计方法的检出功效与错误发现率：本书提出的结构方程模型与整数规划结合的方法（SEMIP）、结构方程模型（SEM）、使用极大似然的结构方程模型（稀疏 MLE；Cai et al. 2013）、QTLnet（Neto et al. 2008），以及通过加权相关网络分析推断的共关联（相关）网络。图2.15给出了这五种统计量随样本量变化的检出功效和错误发现率。可以观察到两个显著特点。第一，SEMIP 优于其余四种方法，其检出功效最高，错误发现率最小。第二，随着样本量增加，五种方法的检出功效均提高，错误发现率均降低。第三，其余四种方法的错误发现率都很高。出乎预料的是，即使样本量为2000，这些方法的错误发现率仍超过60\%。这说明，许多常用统计方法若不通过组合优化搜索具有最佳评分的网络，可能不适合因果网络分析。

\begin{figure}[htbp]
\centering
\begin{minipage}{.49\linewidth}
\centering
\begin{tikzpicture}
\node[inner sep=0,anchor=south west](ch0215imga)at(0,0){\includegraphics[width=\linewidth]{figures/ch02fig15a.pdf}};
\begin{scope}[x={(ch0215imga.south east)},y={(ch0215imga.north west)}]
\fill[white](0,.38)rectangle(.043,.63);
\fill[white](.43,.038)rectangle(.58,.075);
\node[inner sep=0,rotate=90,font=\fontsize{7}{8}\selectfont]at(.021,.51){检出功效};
\node[inner sep=0,font=\fontsize{7}{8}\selectfont]at(.51,.056){样本量};
\fill[white](.125,.675)rectangle(.40,.89);
\draw[black,line width=.3pt](.125,.675)rectangle(.40,.89);
\foreach \yy/\txt in {.86/SEMIP,.825/SEM,.79/{稀疏 MLE},.755/QTLnet,.72/共关联}
 \node[anchor=west,inner sep=0,font=\fontsize{5}{5.7}\selectfont]at(.213,\yy){\txt};
\draw[red!85!black,line width=.65pt](.14,.86)--(.20,.86);
\draw[red!85!black,line width=.65pt,dashed](.14,.825)--(.20,.825);
\draw[blue!65!black,line width=.65pt](.14,.79)--(.20,.79);
\draw[blue!65!black,line width=.65pt,dashed](.14,.755)--(.20,.755);
\draw[green!65!black,line width=.65pt](.14,.72)--(.20,.72);
\end{scope}
\end{tikzpicture}
\end{minipage}\hfill
\begin{minipage}{.49\linewidth}
\centering
\begin{tikzpicture}
\node[inner sep=0,anchor=south west](ch0215imgb)at(0,0){\includegraphics[height=.826\linewidth]{figures/ch02fig15b.pdf}};
\begin{scope}[x={(ch0215imgb.south east)},y={(ch0215imgb.north west)}]
\fill[white](0,.36)rectangle(.043,.64);
\fill[white](.45,.038)rectangle(.60,.075);
\node[inner sep=0,rotate=90,font=\fontsize{7}{8}\selectfont]at(.021,.51){错误发现率};
\node[inner sep=0,font=\fontsize{7}{8}\selectfont]at(.525,.056){样本量};
\fill[white](.17,.285)rectangle(.455,.485);
\draw[black,line width=.3pt](.17,.285)rectangle(.455,.485);
\foreach \yy/\txt in {.455/SEMIP,.42/SEM,.385/{稀疏 MLE},.35/QTLnet,.315/共关联}
 \node[anchor=west,inner sep=0,font=\fontsize{5}{5.7}\selectfont]at(.26,\yy){\txt};
\draw[red!85!black,line width=.65pt](.185,.455)--(.25,.455);
\draw[red!85!black,line width=.65pt,dashed](.185,.42)--(.25,.42);
\draw[blue!65!black,line width=.65pt](.185,.385)--(.25,.385);
\draw[blue!65!black,line width=.65pt,dashed](.185,.35)--(.25,.35);
\draw[green!65!black,line width=.65pt](.185,.315)--(.25,.315);
\end{scope}
\end{tikzpicture}
\end{minipage}

\origfig{15}{用于因果表型网络推断的三种方法的(a)检出功效和(b)错误发现率曲线。}
\end{figure}

\subsection{真实数据分析}

为进一步评价其性能，将所提出的 SEMIP 应用于 UK-2K 数据集。UK2K 队列项目使用低测序深度的全基因组测序（whole-genome sequencing，WGS），评价遗传变异对64种不同性状的贡献（Walter et al. 2015）。不过，许多个体存在表型缺失。为保证纳入个体的表型没有缺失，分析选取765名个体，他们共有33746个基因中的2240049个 SNP，并共同具有13个主要表型组中的39种性状，涵盖了广泛的性状。对表型作基于秩的逆正态变换，并以变换后的值作为性状值，以保证性状值服从正态分布。

为研究包含罕见变异的基因型—表型网络，使用基于基因的函数型结构方程建立基因型—表型网络，作为覆盖真实因果基因型—表型网络的初始网络。推断全基因组范围基因型—表型网络的过程包括四步。第一步，估计每个基因的函数主成分得分。第二步，使用二次正则化函数型典型相关分析（quadratically regularized functional canonical correlation analysis，QRFCCA），识别与39种性状显著关联的全部基因。第三步，使用函数型结构方程模型（FSEM）估计初始网络结构。最后一步，将函数型结构方程模型与整数规划方法结合，推断基因型—表型网络；其中，函数型结构方程用于计算评分函数。

将 QRFCCA 应用于 UK-2K 数据集，共识别出79个与39种性状显著关联、仅含罕见变异的基因。随后，使用39种性状与79个基因构建基因型—表型网络。图2.16显示了最大的连通因果基因型—表型网络，包含39个性状节点、58个基因节点和323条有向边。红色节点表示性状，绿色节点表示基因；黑色实线表示表型之间的因果关系，蓝色虚线表示基因型与表型之间的因果关系（Wang 2016）。

\begin{figure}[htbp]
\centering
\includegraphics[width=.99\linewidth]{figures/ch02fig16.pdf}
\origfig{16}{估计得到的基因型—表型网络，包含39个表型节点和58个基因节点。节点保留原始性状与基因标识符；红色为表型、绿色为基因，黑色实线为表型间的因果关系，蓝色虚线为基因型与表型间的因果关系。}
\end{figure}

变量之间既存在直接因果效应，也存在间接因果效应。直接因果效应指一个变量的变异直接影响另一个变量的变化；间接因果效应指一个变量的变异经由其他变量的中介作用，影响另一个变量的变化。通过路径搜索，在网络中找到1112对因果关系。其中，916对具有间接因果效应，占全部因果关系的82.4\%。

表2.5列出了20对表型—表型及基因型—表型因果关系。直接效应是从原因节点指向结局终点节点的直接边的路径系数。间接效应是从原因节点经过中间中介节点到达结局终点节点的路径上，各路径系数的乘积。总效应为直接效应与间接效应之和。第1章已说明，如果包括直接效应和间接效应在内的所有因果路径均被发现，边际回归效应就等于总效应。如果一个变量的边际效应和总效应作用于另一个变量，那么这两个变量之间的因果路径尚未完全发现。从表2.5可以观察到两个显著特点。第一，所有原因表型与基因型都同时具有直接效应和间接效应。第二，三对表型—表型关系——PLT Alkaline、$\mathrm{WBC}\to\mathrm{FEV1}$ 和 $\mathrm{WBC}\to\mathrm{FVC}$，以及一对基因型—表型关系 $\mathrm{RP11\text{-}363E6.4}\to\mathrm{ApoA1}$——的边际效应与总效应近似相等。不过，其余16个原因表型或基因型的边际效应，均大于或小于相应的总效应。这表明，许多混杂因子没有测量，或未纳入分析。

\begin{table}[htbp]
\centering
\origtable{5}{20对表型—表型及基因型—表型因果关系。响应变量与预测变量保留数据中的标识符。}
\begin{adjustbox}{max width=\linewidth}
\small
\begin{tabular}{llrrrr}
\toprule
响应变量&预测变量&直接效应&间接效应&总效应&边际效应\\\midrule
Alkaline&PLT&0.1397&0.0053&0.1450&0.1467\\
ApoA1&HDL&0.8451&$-0.0369$&0.8082&0.7814\\
ApoA1&VLDL&0.0840&$-0.3284$&$-0.2444$&$-0.1687$\\
ApoB&HDL&$-0.1877$&0.2109&0.0233&$-0.1789$\\
FEV1&WBC&$-0.1577$&$-0.0260$&$-0.1837$&$-0.1818$\\
FVC&VLDL&$-0.0610$&$-0.0131$&$-0.0740$&$-0.1625$\\
FVC&WBC&$-0.1500$&$-0.0309$&$-0.1809$&$-0.1798$\\
Glucose&Insulin&$-0.4846$&0.9631&0.4785&0.4668\\
HDL&TotalCholesterol&2.1633&$-1.9076$&0.2556&0.1939\\
MCH&Hgb&1.1621&0.6690&1.8311&0.2664\\
Total\_Fat\_Mass&VLDL&0.1217&0.0261&0.1477&0.2667\\
Alkaline&AC008694.3&$-0.0901$&0.0181&$-0.0720$&$-0.0634$\\
ApoA1&RP11-363E6.4&0.0534&0.0005&0.0539&0.0553\\
ApoB&RP1-276E15.1&$-0.0915$&$-0.0058$&$-0.0973$&$-0.0723$\\
Bicarbonate&HOXB1&$-0.0950$&0.0112&$-0.0839$&$-0.0709$\\
Bicarbonate&OR4A44P&$-0.0870$&0.0010&$-0.0860$&$-0.0652$\\
Bilirubin&RP11-153M7.3&$-0.0176$&0.0067&$-0.0109$&$-0.0117$\\
Total\_Lean\_Mass&MLN&$-0.0628$&$-0.0056$&$-0.0684$&$-0.0628$\\
Urea&FBXL5&0.0661&0.0008&0.0669&0.0632\\
WBC&RP11-363E6.4&$-0.0946$&0.0842&$-0.0104$&$-0.0133$\\
\bottomrule
\end{tabular}
\end{adjustbox}
\end{table}
\FloatBarrier

\section*{软件包}
\addcontentsline{toc}{section}{软件包}
\markright{软件包}

包含贝叶斯 Dirichlet 等价（BDe）评分与最小描述长度（MDL）评分度量的精确贝叶斯网络学习软件，可以在 G6G 组学与智能软件目录中找到。
\begin{flushleft}\small
\url{http://g6g-softwaredirectory.com/bio/cross-omics/pathway-analysis-grns/20697-Univ-Warsaw-BNFinder.php}
\end{flushleft}
\texttt{bnlearn} 是一个 R 包，其中包含若干用于学习离散变量或连续变量贝叶斯网络结构的算法，同时实现了基于约束和基于评分的算法（\url{https://arxiv.org/pdf/0908.3817}）。

\texttt{bene} 包包含使用可分解评分 AIC、BIC、BDe、fNML 和 LOO，构造全局最优贝叶斯网络结构的软件（\url{https://github.com/tomisilander/bene}）。R 包 \texttt{pcalg}（图模型与因果推断方法）提供了使用图模型进行因果结构学习与因果推断的函数。主要的因果结构学习算法为 PC（用于没有隐藏变量的观测数据）、FCI 与 RFCI（用于存在隐藏变量的观测数据），以及 GIES（用于观测研究数据与干预实验数据混合、且没有隐藏变量的情形）。在因果推断方面，该包实现了 IDA 算法、广义后门准则（GBC）和广义调整准则（GAC）。\texttt{pcalg} 可从
\url{https://cran.r-project.org/web/packages/pcalg/index.html} 下载。

GOBNILP（使用整数线性规划进行全局最优贝叶斯网络学习）是一个 C 程序，可根据完整的离散数据或局部评分学习贝叶斯网络。GOBNILP 使用约束整数规划框架 SCIP。下载地址为
\url{https://www.cs.york.ac.uk/aig/sw/gobnilp/}。

COMICN 包用于下一代测序组学数据的因果推断，能够根据观测数据和干预数据精确学习因果网络。COMICN 采用基于核心的方法学习因果网络，将网络惩罚逻辑回归、结构方程、函数型结构方程和整数规划结合起来。COMICN 可从
\url{http://www.sph.uth.tmc.edu/hgc/faculty/xiong/index.htm} 下载。
